\section{What a rest energy breaks, and the two senses of rotation}
\label{sec:massbreaks}

The arena carries a mass in two readings. In the field-strength reading the rest energy of a field point is a function of the frame
vector, studied in \cref{sec:energy}; the duality rotation $e^{\vartheta G}$ generated by the binding root multiplies the complex
invariant $q$ by $e^{2i\vartheta}$, so for $q\neq0$ only $\vartheta\in\pi\Z$ fixes it (\cref{sec:orth}). In the Dirac reading the mass is
the kernel $T=e_1e_2e_3=\gamma^0J$ of the Grassmann bilinear (\cref{prop:kernels1,cor:fullmass}). This section asks what the second
breaks and compares the rest motion it generates with the binding rotation of the frame. In brief: the Dirac mass breaks no generator
of $\so{3,3}$, and what reduces the six-dimensional symmetry to the interior one is the kinetic term; it breaks the chiral and
Pauli--G\"ursey symmetries of the massless kinetic term down to $\uu1_J$, along three explicit directions; and the two senses of the
Dirac rest rotation are exchanged by charge conjugation, those of the frame's binding rotation by parity. Units are $\hbar=c=1$ unless
a symbol is shown.

\subsection{The Dirac mass: invariant in six dimensions, breaking the kinetic algebra}
\label{sec:massdirac}

With $\gamma^0=e_1$ and $\gamma^a=e_{a+3}$ the kinetic kernels are $K^\mu=-\gamma^0\gamma^\mu$ (\cref{prop:kernels1}). A Grassmann bilinear
$\psi^{\!\top}M\psi$ is invariant under $\delta\psi=Y\psi$ exactly when $Y^{\!\top}M+MY=0$, and by \cref{cor:fullmass} the antisymmetric
kernels with this property for all fifteen bivectors $Y=e_ae_b$ are the multiples of $T$. Write $\mathrm{Spin}(3,3)$ for the group
generated by the exponentials of the bivectors.

\begin{proposition}[$\mathrm{Spin}(3,3)$ keeps the mass form, not the kinetic forms]
\label{mb:Tinv}
On the spinor module $\R^8$:
\begin{enumerate}[label=(\roman*),leftmargin=2.2em]
\item For every bivector $Y=e_ae_b$ and every real $s$, $(e^{sY})^{\!\top}T\,e^{sY}=T$.
\item An element $Y\in\so{3,3}$ satisfies $Y^{\!\top}K^\mu+K^\mu Y=\Lambda^\mu{}_\nu K^\nu$ for some real $4\times4$ matrix $\Lambda$ if and only if
$Y\in\lint\oplus\R J$, a space of dimension seven.
\end{enumerate}
\end{proposition}

\begin{proof}
(i) By \cref{lem:blade} the trivector $T$ commutes with a bivector that shares two or no indices with $\{1,2,3\}$ and anticommutes with a
mixed one, and by \cref{thm:sig} the former are the antisymmetric bivectors and the latter the symmetric ones. For antisymmetric $Y$,
$Y^2=-\id$ and $e^{sY}=\cos s\,\id+\sin s\,Y$ is orthogonal and commutes with $T$. For symmetric $Y$, $Y^2=\id$, $e^{sY}=\cosh s\,\id+\sinh s\,Y$ is
symmetric, and $(e^{sY})^{\!\top}T=e^{sY}T=Te^{-sY}$. (ii) The conditions are $256$ linear equations in the fifteen bivector coefficients
of $Y$ and the sixteen entries of $\Lambda$. Their solution space is seven-dimensional and projects injectively onto the bivectors. It
contains $\lint$, under which the $K^\mu$ transform covariantly, and $J$, which is antisymmetric and commutes with every $\gamma^0\gamma^\mu$,
so that $J^{\!\top}K^\mu+K^\mu J=[K^\mu,J]=0$; these span seven dimensions.
\end{proof}

Item (i) is \cref{cor:fullmass} carried from the Lie algebra to the group; item (ii) is fixed by the algebra. Within $\mathrm{Spin}(3,3)$
the Dirac mass breaks nothing. What reduces the six-dimensional symmetry to $\lint\oplus\R J$ is the kinetic term, in which the clock
$\gamma^0$ is paired with the spatial triple. The breaking that a Dirac mass does describe lives in the internal symmetries of the
kinetic term. Let $\mathfrak g_{\rm massless}=\R\,\Omega\oplus\operatorname{span}\{J,e_2\Omega,e_3\Omega\}\cong\uu1_\Omega\oplus\su2$ be the
algebra of \cref{prop:gauge1}, of which invariance of the mass form leaves $\R J$.

\begin{proposition}[What the Dirac mass breaks]
\label{mb:massless}
\begin{enumerate}[label=(\roman*),leftmargin=2.2em]
\item On $\mathfrak g_{\rm massless}$ the map $Y\mapsto Y^{\!\top}T+TY$ has rank three and image $\operatorname{span}\{T\Omega,\,Te_2\Omega,\,Te_3\Omega\}$,
contained in the six-dimensional kernel space of the interior Lorentz algebra (\cref{thm:subkernels}). The kernel $T\Omega$ commutes
with $J$ and is the pseudoscalar mass of \cref{prop:massangle}; $Te_2\Omega$ and $Te_3\Omega$ anticommute with $J$ and are Majorana kernels.
\item For $M=T(m_s+m_p\Omega)$ the field equation $K^\mu\partial_\mu\Psi=M\Psi$ of $\half\Psi^{\!\top}(K^\mu\partial_\mu-M)\Psi$, with
$\partial_a\Psi=Jp_a\Psi$, is $\dot\Psi=-JH\Psi$, where
\begin{equation}
\label{mb:eq:Hmp}
H=\gamma^0\gamma^ap_a+\gamma^0(m_s+m_p\Omega),\qquad H^{\!\top}=H,\qquad [H,J]=0,\qquad H^2=\bigl(|\bp|^2+m_s^2+m_p^2\bigr)\id .
\end{equation}
\item The chiral rotation $e^{\xi\Omega}$ preserves every $K^\mu$ and sends $M$ to $(e^{\xi\Omega})^{\!\top}Me^{\xi\Omega}=Me^{2\xi\Omega}$; the stabiliser
of $M\neq0$ in $U(1)_\Omega$ is $\{\pm\id\}$.
\end{enumerate}
\end{proposition}

\begin{proof}
(i) Every element of $\mathfrak g_{\rm massless}$ is antisymmetric, so the map is $Y\mapsto[T,Y]$. By \eqref{eq:bladecomm} $T$ anticommutes
with $\Omega$, $e_2\Omega$ and $e_3\Omega$ and commutes with $J$, which gives the images $2T\Omega$, $2Te_2\Omega$, $2Te_3\Omega$ and $0$. They are
Lorentz-invariant kernels because $T$ is one and $\mathfrak g_{\rm massless}$ commutes with $\lint$. $T$ and $\Omega$ commute with $J$, while
$e_2$ and $e_3$ anticommute with it. (ii) Varying gives $K^\mu\partial_\mu\Psi=M\Psi$, since $K^\mu$ is symmetric and $M$ antisymmetric. With
$K^0=-\id$, and $J$ commuting with $\gamma^0\gamma^a$, $\gamma^0$ and $\Omega$, $\dot\Psi=(K^aJp_a-M)\Psi=-JH\Psi$. The five matrices
$\gamma^0\gamma^a$, $\gamma^0$ and $\gamma^0\Omega=e_4e_5e_6$ are symmetric, square to $\id$ and anticommute pairwise, which gives $H^2$.
(iii) $\Omega\in\mathfrak g_{\rm massless}$ is antisymmetric and anticommutes with $T$, so $e^{-\xi\Omega}Te^{\xi\Omega}=Te^{2\xi\Omega}$; and
since $M^{\!\top}M=(m_s^2+m_p^2)\id$, a nonzero $M$ is invertible, so $Me^{2\xi\Omega}=M$ exactly when $e^{2\xi\Omega}=\id$, that is for $\xi\in\pi\Z$, where $e^{\xi\Omega}=\pm\id$.
\end{proof}

The count of three in (i) is that of \cref{thm:twophases}(a); the explicit broken directions are fixed by the algebra. Their physical
content, that a Dirac mass breaks the axial $U(1)$ and the Pauli--G\"ursey $SU(2)$ of a massless Dirac field down to fermion number
\cite{Pauli57,Gursey58}, is inherited through \cref{thm:realification}. Items (ii) and (iii) are inherited: they are the realification of
$h=\boldsymbol\alpha\cdot\bp+\beta(m_s+i\gamma_5m_p)$ and of the chiral rotation, here derived from the Grassmann kernels; the commuting
first-order form of \cref{sec:dirac} with the same $H$ has the same flow. The two readings share a shape, visible already in the bilinears of \cref{tab:diracmaxwell}. In both, the mass is the
modulus of an order parameter of charge two under a $U(1)$ generated by a complex structure commuting with a Lorentz algebra, $G$ for
$\lG$ and $\Omega$ for $\lint$, and a nonzero modulus leaves only $\{\pm\id\}$. They differ in kind: the Dirac mass is a coupling, or a value
selected by a potential (\cref{thm:twophases}), while the rest energy of a frame vector is a property of a configuration that no
potential selects, the minimum of the energy density being the zero field \eqref{eq:EBbound}.

\subsection{Particle and antiparticle: the two senses of one rotation}
\label{sec:masssenses}

At $\bp=0$ the Hamiltonian \eqref{eq:Hp} of the commuting first-order form $\cL_1=\half\Psi^{\!\top}J\dot\Psi-\half\Psi^{\!\top}H\Psi$ is
$H(0)=m\gamma^0=me_1$.

\begin{proposition}[The rest flow]
\label{mb:restflow}
$JH(0)=mT$, so the rest flow is $\Psi(\lambda)=e^{-m\lambda T}\Psi_0=(\cos m\lambda\,\id-\sin m\lambda\,T)\Psi_0$. On the eigenspaces $e_1=\pm1$, each
four-dimensional and $J$-invariant, $T=\pm J$, so $\Psi(\lambda)=e^{\mp m\lambda J}\Psi_0$ with energy $\half\Psi^{\!\top}H\Psi=\pm\half m|\Psi|^2$.
Moreover $\gamma^0=-TJ$.
\end{proposition}

\begin{proof}
$Je_1=e_2e_3e_1=e_1e_2e_3=T$. Since $T=e_1J$ with $[e_1,J]=0$, $T\Pi_\pm=\pm J\Pi_\pm$ for $\Pi_\pm=\half(\id\pm e_1)$; $T^2=-\id$ gives the
exponential, and $-TJ=-e_1J^2=e_1$.
\end{proof}

The proposition is inherited: it is the realification of $\psi=e^{\mp im\lambda}\psi_0$ on $\beta=\pm1$, at angular frequency $mc^2/\hbar$ in
restored units. The arena adds a name. The rest motion is the one-parameter group of the temporal volume element $T$, the unique
$\mathrm{Spin}(3,3)$-invariant mass form (\cref{cor:fullmass,mb:Tinv}), and the clock $\gamma^0=-TJ$ is the quotient of that rotation by the charge
rotation. Read in this way, the single rotation $e^{-m\lambda T}$ splits into the two senses of the charge rotation according to the
eigenvalue of the clock.

The frame has a rotation with two senses as well. With the split-form Lagrangian $\cL_\kappa=\tfrac m2\eta(\alpha_+X,\alpha_+X)$ of
\eqref{eq:splitL}, every solution is $z(\tau)=c_-e^{-i\kappa\tau}+c_+e^{i\kappa\tau}$ with $c_\pm\in\C^3$ (\cref{prop:orbitsEM}). Below
$\mathsf C$, $\mathsf P$, $\mathsf T$ name the operations. On spinors they are realised by the operators of \cref{thm:C,thm:PT}; on the frame
by $C_f=-\id$, $P_f=\eta$ ($\bE\mapsto-\bE$) and $T_f\colon X(\tau)\mapsto-\eta X(-\tau)$, which combines the reversal of $\tau$ with the field
map $(\bB,\bE)\mapsto(-\bB,\bE)$ of $C_fP_f$, as $T=CP$ on fields in \cref{thm:achiral}.

\begin{proposition}[Which operation reverses the sense]
\label{mb:senses}
\begin{enumerate}[label=(\roman*),leftmargin=2.2em]
\item Every unit vector $C$ of the surplus plane $U$ satisfies $C^2=\id$, $[C,T]=0$, $\{C,J\}=\{C,e_1\}=0$ and $C^{\!\top}H(0)C=-H(0)$. It maps the
solution $e^{-m\lambda J}\Psi_0$ on $e_1=+1$ to the solution $e^{+m\lambda J}C\Psi_0$ on $e_1=-1$ of the same flow. The parities $\gamma^0$ and
$\gamma^0J$ and the time reversal $T_{\rm rev}=e_2e_4e_5e_6$ of \cref{thm:PT} keep the sign of the rest energy.
\item On the frame $P_fGP_f^{-1}=-G$, and on $z=c_-e^{-i\kappa\tau}+c_+e^{i\kappa\tau}$
\[
P_f\colon(c_-,c_+)\mapsto(\bar c_+,\bar c_-),\qquad C_f\colon(c_-,c_+)\mapsto(-c_-,-c_+),\qquad T_f\colon(c_-,c_+)\mapsto(-\bar c_-,-\bar c_+).
\]
\item The sense of the Dirac rest rotation is reversed by $\mathsf C$, $\mathsf{CP}$, $\mathsf{CT}$ and $\mathsf{CPT}$ and kept by $\mathsf P$,
$\mathsf T$ and $\mathsf{PT}$. The sense of the frame's binding rotation is reversed by $\mathsf P$, $\mathsf{CP}$, $\mathsf{PT}$ and $\mathsf{CPT}$
and kept by $\mathsf C$, $\mathsf T$ and $\mathsf{CT}$.
\end{enumerate}
\end{proposition}

\begin{proof}
(i) Write $C=\cos\sigma\,e_2+\sin\sigma\,e_3$. By \eqref{eq:bladecomm} it anticommutes with $e_1$, with $J$ and with the $\gamma^a$, and
commutes with $T$; it is symmetric with $C^2=\id$. Hence $C^{\!\top}e_1C=-e_1$, $C$ maps the eigenspace $e_1=+1$ onto $e_1=-1$, and $C$ commutes
with $e^{-m\lambda T}$. The operators $\gamma^0=e_1$ and $\gamma^0J=T$ commute with $e_1$ and with $T$, so they preserve each eigenspace and the
flow. $T_{\rm rev}$ commutes with $e_1$ and anticommutes with $J$ and $T$, so $T_{\rm rev}\Psi(-\lambda)=e^{-m\lambda T}T_{\rm rev}\Psi_0$ solves the
same flow on the same eigenspace. (ii) $\eta G\eta=G^{\!\top}=-G$. On $z$, $P_f$ acts as $z\mapsto\bar z$, $C_f$ as $z\mapsto-z$, and $T_f$ as
$z(\tau)\mapsto-\overline{z(-\tau)}$; substituting the general solution gives the three maps. (iii) Each operation either exchanges the two
senses or keeps both, so the signs multiply under composition; the single operations are (i), (ii) and \cref{thm:PT}.
\end{proof}

Item (i) is inherited, realised by the arena's charge conjugation and time reversal (\cref{thm:C,thm:PT}). Items (ii) and (iii) are
fixed by the algebra; (ii) makes explicit, on the two senses, the action of the operations of \cref{thm:EBsplit}, by which $P_f$ maps
$\cL_\kappa$ to $\cL_{-\kappa}=\cL_\kappa+2m\kappa\,\tfrac{d}{d\tau}(\bE\cdot\bB)$, a symmetry up to a total derivative, while $C_f$ and $T_f$ leave $\cL_\kappa$ invariant.

For Dirac matter the two senses are particle and antiparticle. They are solutions of one Lagrangian, of angular frequency $m$, each with
energy of definite sign, and they are exchanged by $C$, a symmetry of the Grassmann Lagrangian, since it preserves the mass form and every
$K^\mu$ (\cref{thm:C}); on the commuting first-order form $\cL_1$ it reverses the sign of the action, because $CJC^{-1}=-J$ and
$CH(\bp)C^{-1}=-H(-\bp)$. For the frame the two senses are the kernels of the two factors of $\alpha_-\alpha_+=\partial_\tau^2+\kappa^2$. One
carries no momentum and the other all of it, and each alone carries zero energy, since the energy is $2m\kappa^2\Re(c_-\cdot c_+)$
(\cref{prop:orbitsEM}). They are exchanged by parity, which is a symmetry only up to a total derivative, and their frequency
$\kappa=c/\ell$ \eqref{eq:kappa} is a conversion rate, not a mass. In the interior $m$ and $\kappa$ are independent inputs: $m$ multiplies the
Lagrangian and $\kappa$ sets the period. Reading the frame's $\hbar\kappa$ as $mc^2$ would therefore be an added identification, not a result.
That the frame's binding is not a particle--antiparticle pair in the Dirac sense is our reading of these facts. In both cases the
reversal is $\mathsf{CPT}$-odd.

\subsection{Two vacua}
\label{sec:masstwovacua}

The Fock vacuum of one Dirac momentum is unique and neutral, with energy $-2E$, and charge conjugation lifts to a unitary operator on
the Fock space that anticommutes with the charge and carries $H(\bp)$ to $H(-\bp)$ (\cref{thm:spectra,prop:Cstates}). At rest its energy,
$-2m$ ($-2mc^2$ in restored units), is minus half the sum of the rest energies of the four modes. Only in that sense does the vacuum
carry the rest energy with the opposite sign, and only this vacuum does.

The frame has a different vacuum. On a coordinate subspace with $k$ temporal and $l$ spatial directions the zero-point energy of the
interior oscillators is $\tfrac12\hbar\kappa(k-l)$, and it vanishes on $\R^{3,3}$ (\cref{thm:subfs}). The full frame has no state of
lowest energy for this to be the energy of.

\begin{corollary}[No ground state]
\label{mb:noground}
Quantise the interior oscillators of the full frame canonically, mode by mode, with positive norm. Then the spectrum of the energy
$\mathcal E=\tfrac m2\bigl(\eta(\dot X,\dot X)+\kappa^2\eta(X,X)\bigr)$ is $\hbar\kappa\Z$, every level is infinitely degenerate, and every state has one
of lower energy. Making the spatial energies positive instead requires negative norm.
\end{corollary}

\begin{proof}
Since $G$ is an anti-isometry, $\mathcal E$ is the energy of \eqref{eq:splitEg}, and
$\mathcal E=\sum_a\eta_{aa}\,\tfrac m2(\dot x_a^2+\kappa^2x_a^2)$. The cross term of $\cL_\kappa$ is a total derivative (\cref{thm:EBsplit}); it
shifts the momenta by a canonical transformation and leaves the spectrum unchanged. A spatial mode $x$ has Lagrangian
$-\tfrac m2(\dot x^2-\kappa^2x^2)$, momentum $p_x=-m\dot x$ and Hamiltonian $-(p_x^2/2m+m\kappa^2x^2/2)$. With
$x=\sqrt{\hbar/2m\kappa}\,(\hat a+\hat a^\dagger)$ and $p_x=i\sqrt{\hbar m\kappa/2}\,(\hat a^\dagger-\hat a)$,
$p_x^2/2m+m\kappa^2x^2/2=\tfrac{\hbar\kappa}2(\hat a\hat a^\dagger+\hat a^\dagger\hat a)$, so a temporal mode contributes $\hbar\kappa(j+\half)$ and a spatial one
$-\hbar\kappa(j+\half)$, $j=0,1,\dots$. On $\R^{3,3}$ the half-integers cancel and a state of occupations $(j_1,j_2,j_3;j_4,j_5,j_6)$ has energy
$\hbar\kappa(j_1+j_2+j_3-j_4-j_5-j_6)$. Every integer occurs; adding one temporal and one spatial quantum keeps the energy, and adding a
spatial quantum lowers it by $\hbar\kappa$. A positive spatial energy needs $[\hat a,\hat a^\dagger]=-1$ on the spatial modes, a negative norm.
\end{proof}

The zero-point values restate \cref{thm:subfs}, and the per-mode quantisation is the inherited indefinite oscillator. The cancellation is
fixed by the algebra: it occurs because the signature is neutral, which is the condition under which factorisation and binding coexist
(\cref{thm:neutral}). \Cref{mb:noground} is an obstruction under the hypothesis of positive norm, the familiar fate of indefinite
oscillators \cite{PaisUhlenbeck50}, and makes explicit for the full space what \cref{sec:subfields} records for the mixed planes. The two
vacua differ in kind: the Dirac vacuum has energy $-2E$ per momentum and is bounded below, while the frame's zero-point energy cancels
between positive temporal and negative spatial oscillators, and there is no ground state for it to be the energy of.

\subsection{Goldstone directions across the tower}
\label{sec:massgoldstone}

\Cref{mb:tab:goldstone} collects the broken directions at the stages where the paper computes them, at $n$ constituents, counted for the specified generators:
at one constituent the four of $\mathfrak g_{\rm massless}$, at three constituents the $21$ Pati--Salam generators and the central $J$ and
$J\Gamma$ (\cref{prop:global}). One element runs through all of them.

\begin{table}[t]
\centering
\small
\begin{tabular}{@{}>{\raggedright\arraybackslash}p{2.4cm}>{\raggedright\arraybackslash}p{2.3cm}>{\raggedright\arraybackslash}p{1.7cm}>{\raggedright\arraybackslash}p{4.4cm}>{\raggedright\arraybackslash}p{3.0cm}@{}}
\toprule
stage & order parameter & kept & broken directions & source\\
\midrule
$n=1$, Dirac & $T=\gamma^0J$ & $\uu1_J$ & $3$, global: $T\Omega$, $Te_2\Omega$, $Te_3\Omega$ & \cref{thm:twophases}(a), \cref{mb:massless}\\
$n=1$, Majorana & kernel on $\Gamma=+1$ & $\R(J-J\Gamma)$ & $3$, global & \cref{thm:twophases}(b)\\
$n=3$, Pati--Salam & $\Phi_0$ ($\nu_R\nu_R$) & $14$ of $23$ & $9$, eaten by $Z'$, $W_R^\pm$ and six $X$ & \cref{thm:nomajoron}\\
$n=3$, electroweak & $\Phi_0$ and bidoublet & $10$ of $23$ & $12$ eaten, $1$ axion-like ($1-\Gamma$) & \cref{thm:axion}\\
\bottomrule
\end{tabular}
\caption{Broken directions along the tower.}
\label{mb:tab:goldstone}
\end{table}

\begin{lemma}[One charge across the scales]
\label{mb:oneminusGamma}
Let $J$ be the complex structure and $\Gamma$ the chirality, at one constituent or at three, so that $[J,\Gamma]=0$, and let $J-J\Gamma$ be the generator of charge $1-\Gamma$.
\begin{enumerate}[label=(\roman*),leftmargin=2.2em]
\item If a kernel $M\neq0$ commutes with $J$ and anticommutes with $\Gamma$ (Dirac type), then $[J-J\Gamma,M]=-2J\Gamma M\neq0$.
\item If $M$ lives on $\Gamma=+1$, that is $\Gamma M=M\Gamma=M$, and anticommutes with $J$, then $[J-J\Gamma,M]=0$.
\item At one constituent $J-J\Gamma=J+\Omega$.
\end{enumerate}
\end{lemma}

\begin{proof}
$J$ commutes with $\Gamma$. (i) $[J,M]=0$ and $MJ\Gamma=-J\Gamma M$, so $[J\Gamma,M]=2J\Gamma M$, which is nonzero because $J\Gamma$ is invertible.
(ii) $[J,M]=2JM$, and $[J\Gamma,M]=JM-MJ\Gamma=JM+JM\Gamma=2JM$. (iii) $\Gamma=J\Omega$, so $J\Gamma=J^2\Omega=-\Omega$.
\end{proof}

The lemma is inherited: a charge that vanishes on one chirality is kept by a Majorana mass of that chirality and broken by every mass
that pairs opposite chiralities. What the arena adds is the identification across scales: one element of the algebra, $J-J\Gamma$, is
kept by a Majorana mass on $\Gamma=+1$ and broken by every Dirac-type mass, both at one constituent and at three. At one constituent it is
$J+\Omega$; the Majorana phase keeps it, and the Dirac mass $T$, which commutes with $J$ and anticommutes with $\Gamma$, breaks it
(\cref{thm:twophases}). At three constituents $\Phi_0$ lives on $\Gamma=+1$, since $\Phi_0^{\!\top}\Phi_0=P_\nu$ and $P_\nu$ projects into the eigenspace $\Gamma=+1$, and anticommutes with $J$
(\cref{thm:kernelspace}); item (ii) is therefore the one-line reason for the equality $[J,\Phi_0]=[J\Gamma,\Phi_0]$ in \cref{thm:nomajoron}. The bidoublet kernels are of Dirac type (\cref{thm:kernelspace}), so every nonzero bidoublet value breaks
$J-J\Gamma$ by item (i); at the generic electroweak vacuum, with both bidoublet values nonzero, its direction escapes the gauge orbit (\cref{thm:axion}).

\subsection{What a mass breaks}
\label{sec:massverdict}

\begin{table}[t]
\centering
\small
\begin{tabular}{@{}>{\raggedright\arraybackslash}p{3.0cm}>{\raggedright\arraybackslash}p{6.6cm}>{\raggedright\arraybackslash}p{4.6cm}@{}}
\toprule
symmetry & broken by a mass? & where\\
\midrule
interior Lorentz $\lint$ & no: $P^2$ and $T$ are invariant; the little group of a momentum is three-dimensional with or without mass & inherited
\cite{WeinbergQFT}; \cref{sec:energy}\\
frame Lorentz $\lG$ & no: $q$ is invariant; the stabiliser of a nonzero field point is two-dimensional with or without $q$ & inherited
\cite{PenroseRindler84}; \cref{sec:energy}\\
$\mathrm{Spin}(3,3)$ & no: $T$ is invariant; the kinetic term keeps seven dimensions & fixed; \cref{cor:fullmass}, \cref{mb:Tinv}\\
duality $U(1)_G$ & yes, to $\{\pm\id\}$ on Lorentz orbits, exactly when $q\neq0$; for a configuration, not a vacuum & inherited; \cref{sec:orth,sec:energy}\\
axial $U(1)_\Omega$ & yes, to $\{\pm\id\}$ & inherited; \cref{prop:massangle,prop:gauge1}, \cref{mb:massless}\\
Pauli--G\"ursey $\su2$ & yes, to $\uu1_J$ & inherited; \cref{prop:gauge1,thm:twophases}\\
$U(1)_{1-\Gamma}$ & by every Dirac-type mass; not by a Majorana mass on $\Gamma=+1$ & inherited; \cref{thm:twophases,thm:axion}, \cref{mb:oneminusGamma}\\
\bottomrule
\end{tabular}
\caption{What a mass breaks in the arena.}
\label{mb:tab:verdict}
\end{table}

\Cref{mb:tab:verdict} collects the answer. The statement that a mass is what breaks a symmetry is true in the arena for duality in the
frame, in the orbit-space sense and for a configuration rather than a vacuum, and for the chiral and Pauli--G\"ursey symmetries of Dirac
matter. It is false for every spacetime symmetry, the six-dimensional one included: a mass changes the type of a Lorentz stabiliser, never its
dimension, and at a field point what a nonzero $q$ removes is the duality generator (\cref{sec:energy}).

Which term counts as the symmetric part, in the sense of the clause on $c^2$ discussed at the end of \cref{sec:energy}, and which as the breaking depends on which symmetry is meant. For $\mathrm{Spin}(3,3)$ the mass form is invariant and the
kinetic term breaks the symmetry (\cref{mb:Tinv}); for the massless kinetic algebra $\uu1_\Omega\oplus\su2$ it is the mass that cuts it to
$\uu1_J$ (\cref{mb:massless}); and in the frame the generator of duality is the binding root, which enters the interior through the shift
$\kappa G$ but is not an infinitesimal isometry of the split form on which the interior Lagrangian is built (\cref{sec:orth,sec:energy}).
In the closure of the force law the breaking of duality is explicit: with Lorentz-class inertia the duality angle of the momentum field
moves in a potential whose amplitude vanishes exactly when that field is null, that is when its rest energy vanishes
(\cref{thm:pendulum,cor:constrate}). This restates \cref{sec:energy}; the breaking is explicit, carried by the inertia and not by a vacuum, and the reading of the angle as a
pseudo-Goldstone boson of duality, with its limits, is taken up in \cref{sec:restclosure}.
